\paragraph{Practical implications.}
The protocol reveals that communicative function, not message content, determines pruning safety. This finding enables differentiated deployment strategies. Models with terse, content-free acknowledgments (Qwen pattern) tolerate aggressive pruning with minimal risk. Models with verbose, error-carrying acknowledgments (Llama pattern) benefit from pruning precisely because removal eliminates an error-injection pathway. Message suppression rates translate directly to proportional savings in inference calls, latency, and cost. Backbone choice determines the nature of the benefit: pruning is purely a cost optimization for Qwen but doubles as an accuracy improvement for Llama. Practitioners should profile acknowledgment surface form before choosing a suppression threshold; annotation-free heuristic detectors suffice where false positives carry no accuracy risk (Qwen), but higher-precision classifiers are warranted for error-channel backbones. Multi-agent frameworks could also suppress acknowledgments at the source by instructing agents not to produce purely confirmatory replies. The effectiveness of this approach remains uncertain, however, because our results indicate acknowledgments arise as emergent behavior rather than from explicit prompt design. Post-hoc pruning offers a more practical path for deployed systems, as it requires no modification to agent prompts or underlying model behavior. Qwen pruning remains robust across decoding temperatures; Llama heuristic pruning shows temperature sensitivity consistent with its error-channel mechanism (Appendix~\ref{app:temp-robustness}). The protocol deploys to new backbones with no task-specific training data. We conjecture that aggregate removal dissolves self-reinforcing acknowledgment chains (Appendix~\ref{app:cascade}), though prospective validation is needed.

\paragraph{Why do mechanisms differ?}
The divergence likely traces to alignment training. Qwen's instruction tuning produces short, formulaic confirmations (median 4 tokens) that serve a purely structural grounding role~\cite{clark_contributing_1989}. Llama's RLHF training encourages verbose, elaborative responses even for simple confirmations (mean 59.8 tokens), and creates a channel through which fabricated details propagate. In information-theoretic terms, Qwen acknowledgments occupy zero bits of the task-relevant channel, while Llama acknowledgments consume channel capacity with noise that degrades the signal~\cite{shannon_1948}. This pattern echoes the observation that backchannels in human dialogue are load-bearing for grounding but not for propositional content~\cite{bavelas_co_narrators_2000}; whether LLM acknowledgments preserve this separation depends on training incentives. The analogy is imperfect, however. Human backchannels are listener-produced signals, typically nonverbal or minimally verbal, whereas LLM acknowledgments constitute full speaker turns generated by the responding agent. Qwen acknowledgments approximate classic backchannels in brevity and vacuous content, but Llama acknowledgments diverge from the backchannel prototype by producing lengthy turns with substantive (and often inaccurate) elaboration, a pattern driven by RLHF's preference for verbose responses. Acknowledgment behavior is emergent and persists without prompt instructions; the pruning principle generalizes even where heuristic detectors have limited transferability (Appendix~\ref{app:boundary_conditions}).

\paragraph{Limitations.}
Key limitations include single-seed causal evidence (cross-seed replication of injection experiments remains future work) and domain scope limited to mathematical reasoning, commonsense QA, and science QA; generalization to dialogue or code generation tasks is untested.

We discuss limitations including topology scope, decoding sensitivity, and single-seed injection constraints in Appendix~\ref{app:extended_limitations}.
